\documentclass[11pt, a4paper]{article}

\usepackage[T1]{fontenc}
\usepackage[utf8]{inputenc}
\usepackage{tgpagella}
\usepackage{tgcursor} % Fuente monoespaciada robusta
\usepackage{microtype}
\usepackage[spanish, es-noshorthands]{babel}
\usepackage[table, xcdraw]{xcolor} % Soluciona los errores de \color no definido
\usepackage[margin=2.5cm]{geometry}
\usepackage{fancyhdr}
\usepackage{listings}

\lstset{
    language=Python,
    basicstyle=\ttfamily\small,
    backgroundcolor=\color{gray!5},
    frame=single,
    rulecolor=\color{gray!40},
    commentstyle=\color{green!60!black},
    showstringspaces=false,
    breaklines=true,          % Soluciona desbordamientos horizontales
    breakatwhitespace=true,
    columns=flexible
}

\pagestyle{fancy}
\fancyhf{}
\setlength{\headheight}{16pt}
\lhead{\textbf{IMT-342}}
\chead{Notebook Skeleton Hito 2}
\rhead{Semana 9}

\begin{document}

\section*{Esqueleto Oficial de Jupyter Notebook (Hito 2)[cite: 13]}
Estructure su entrega computacional siguiendo esta jerarquía de celdas Markdown y Code.

\begin{lstlisting}
# 1. Planteamiento del Problema y Modelo del Robot
# [Descripcion del modelo, tabla D-H, convenciones]

# 2. Cinematica Inversa Analitica (Metodo de Pieper)
# [Calculo del centro de muneca y orientacion residual]

# 3. Verificacion FK
# [Funcion que evalua candidatos: T_0^6(q) approx T_d]

# 4. Jacobiano Geometrico
# [Derivacion de columnas e implementacion en codigo]

# 5. Analisis de Singularidades
# [Calculo de rangos, determinantes y experimento]

# 6. Generacion de Trayectoria (LSPB)
# [Funciones q(t), qdot(t), qddot(t) y limites]

# 7. Dinamica / Analisis de Torque
# [PLACEHOLDER: A desarrollar despues de W09-S02]

# 8. Graficas y Discusion
# [Visualizacion del perfil LSPB y torques]

# 9. Validacion y Conclusiones
# [Argumentacion sobre la validez del algoritmo]
\end{lstlisting}

\end{document}
